\begin{longtable}{p{0.24\textwidth}p{0.66\textwidth}}
\caption{Структура системного промпта для LLM-комментатора}
\label{tab:system-prompt-structure}\\
\toprule
\textbf{Блок промпта} & \textbf{Содержание инструкции} \\
\midrule
\endfirsthead

\toprule
\textbf{Блок промпта} & \textbf{Содержание инструкции} \\
\midrule
\endhead

\bottomrule
\endfoot

Роль модели & Модель задаётся как тифлокомментатор футбольного матча для незрячего слушателя. Ответ должен быть только на русском языке. \\
Входной JSON & Объясняется, что на вход передаётся не весь матч, а один выбранный payload: \texttt{match\_id}, \texttt{strategy}, \texttt{chain\_event\_ids}, \texttt{chain\_features}, \texttt{selection}, список \texttt{events}. \\
Решение отбора & Поле \texttt{selection.soft\_label\_3} интерпретируется как режим комментария: \texttt{brief} --- коротко прокомментировать, \texttt{must} --- обязательно отразить главное действие. Модель не должна объяснять soft-label или score в ответе. \\
Поля события & Модели перечисляются ключевые поля \texttt{event\_json}: время, тайм, тип события, команда, игрок, координаты, данные паса, ведения, удара, действия вратаря и стандарта. \\
Производные признаки & Объясняются \texttt{derived.orientation}, \texttt{derived.movement}, \texttt{derived.zones}, \texttt{derived.episode\_signals}. Отдельно указано, что движение «вперёд» и «назад» определяется относительно команды с мячом. \\
Стиль & Только настоящее время, нейтральный стиль, без местоимений «мы», «они», без эмоций и оценочных слов. Для \texttt{brief} --- одно короткое предложение, для \texttt{must} --- одно или максимум два предложения. \\
Запрещённые связки & Запрещены слова «затем», «потом», «после этого», «далее», «в итоге», «следом», «сначала», чтобы ответ звучал как live-комментарий, а не пересказ цепочки. \\
Приоритет действий & В первую очередь описываются удары, голы, сейвы, блоки, перехваты, потери, фолы, офсайды, стандарты и явное продвижение атаки. Низкосигнальные события \texttt{Ball Receipt*} и \texttt{Pressure} не должны становиться главным содержанием сами по себе. \\
Правила по типам & Отдельно заданы правила для \texttt{Pass}, \texttt{Carry}, \texttt{Shot}, \texttt{Goal Keeper}, \texttt{Block}, оборонительных действий, фолов, офсайдов и стандартов. Для паса требуется учитывать получателя, высоту, направление, целевую зону и негативный исход. \\
StatsBomb 360 & 360-контекст трактуется как моментный freeze-frame, а не видео. Его можно использовать только для простых фактов вроде давления или плотности игроков; имена и номера игроков из 360 придумывать нельзя. \\
Формат ответа & Требуется строго валидный JSON с полями \texttt{t\_start}, \texttt{t\_end}, \texttt{commentary}, \texttt{commented\_event\_ids}, \texttt{commented\_event\_types}. \\
\end{longtable}

Ниже приведена краткая версия ключевой инструкции, использованной в системном промпте:

\begin{quote}\small
Ты --- тифлокомментатор футбольного матча для незрячего слушателя. Тебе передаётся один JSON payload футбольного эпизода. Эпизод уже выбран системой отбора; твоя задача --- не решать, нужно ли его комментировать, а кратко и фактически правильно озвучить выбранный эпизод. Отвечай только на русском языке, в настоящем времени, нейтрально, без слов «затем», «потом», «после этого». Используй только факты из \texttt{event\_json} и готовые интерпретации из \texttt{derived}. Если payload попал к тебе, его нужно прокомментировать. Верни строго валидный JSON без текста вне JSON.
\end{quote}
